%% ma: L12-B18-0002
%% lop: 12
%% chuong: 6
%% bai: 18
%% dang: 12.18.2
%% loai: DS
%% muc: [TH, TH, TH, TH]
%% dap_an: "SSĐĐ"
%% nguon: "(NBV)-ĐỀ SỐ 5-MỖI NGÀY 1 ĐỀ THI 2025.pdf (D05, MỖI NGÀY 1 ĐỀ - Đề số 5), Phần II Câu 3"
%% kien_thuc: "Bài 18 (xác suất có điều kiện)"
%% trang_thai: cho-duyet
%% ly_do: "Dạng toán mới đề xuất 12.18.2: xác suất có điều kiện khi lấy không hoàn lại (Đúng/Sai)"
%% ghi_chu: "Bỏ ảnh két nước ngọt: ảnh minh họa đời sống, không chứa dữ kiện toán"
\begin{ex}
Một két nước ngọt đựng $24$ chai nước có khối lượng và hình thức bề ngoài như nhau, trong đó có $16$ chai loại I và $8$ chai loại II. Bác Tùng lần lượt lấy ngẫu nhiên hai chai nước (lấy không hoàn lại). Xét các biến cố: $A$: ``Lần thứ nhất lấy ra chai nước loại I''; $B$: ``Lần thứ hai lấy ra chai nước loại I''.
\choiceTF
{$P(B\mid A)=\dfrac{16}{23}$.}
{$P(B\mid\overline{A})=\dfrac{15}{23}$.}
{\True $P(\overline{B}\mid A)=\dfrac{8}{23}$.}
{\True $P(\overline{B}\mid\overline{A})=\dfrac{7}{23}$.}
\loigiai{a) Nếu lần đầu lấy chai loại I thì còn $23$ chai, trong đó $15$ chai loại I nên $P(B\mid A)=\dfrac{15}{23}$. Sai.
b) Nếu lần đầu lấy chai loại II thì còn $23$ chai, trong đó $16$ chai loại I nên $P(B\mid\overline{A})=\dfrac{16}{23}$. Sai.
c) $P(\overline{B}\mid A)=1-\dfrac{15}{23}=\dfrac{8}{23}$. Đúng.
d) $P(\overline{B}\mid\overline{A})=1-\dfrac{16}{23}=\dfrac{7}{23}$. Đúng.}
\end{ex}
